% Generated draft: requires mathematical and source review.
\documentclass[11pt]{article}
\usepackage[margin=25mm]{geometry}
\usepackage{fontspec}
\setmainfont{Latin Modern Roman}
\usepackage{amsmath,amssymb,mathtools,unicode-math}
\setmathfont{Latin Modern Math}
\usepackage{enumitem,needspace,xurl,hyperref}
\setlength{\parindent}{0pt}
\setlength{\parskip}{5pt}
\setlength{\emergencystretch}{3em}
\newcommand{\N}{\mathbb{N}}
\newcommand{\Z}{\mathbb{Z}}
\newcommand{\Q}{\mathbb{Q}}
\newcommand{\R}{\mathbb{R}}
\newcommand{\C}{\mathbb{C}}
\newcommand{\F}{\mathbb{F}}
\newcommand{\A}{\mathbb{A}}
\newcommand{\PP}{\mathbb{P}}
\newcommand{\E}{\mathbb{E}}
\newcommand{\eps}{\varepsilon}
\newcommand{\dd}{\,\mathrm{d}}
\begin{document}
\section*{UnsolvedMath ID: 30002977}
\textbf{Source problem code:} catalogue entry
\subsection{Definitions and mathematical statement}
Under the usual Severi-variety interpretation, let \((X,L)\) be a smooth projective polarized \(K3\) surface: \(\omega_X\simeq\mathcal O_X\), \(H^1(X,\mathcal O_X)=0\), and \(L\) is ample. The supplied record does not specify the base field, whether \(L\) is primitive, or the polarization moduli space relative to which \((X,L)\) is general; these remain unresolved hypotheses.

The conventional notation is as follows. Write \(L^2=2p-2\). For a positive integer \(n\) and a nonnegative integer \(\delta\), let \(V_{n,\delta}(X)\) denote the locus in \(|L^{\otimes n}|\) parametrizing integral curves with exactly \(\delta\) ordinary nodes and no other singularities. Their normalizations have genus
\[
g=1+n^2(p-1)-\delta.
\]
These conventions explain the notation but are not explicitly supplied in the record and require confirmation against the report. In particular, the intended admissible ranges of \(p,n,\delta\), including any restrictions on \(g\), are not specified.

For a nonempty irreducible component \(V\subset V_{n,\delta}(X)\), the normalization moduli map, under the usual coarse-moduli interpretation of \(M_g\), is
\[
\psi_{X,n,\delta}|_V:V\longrightarrow M_g,
\qquad C\longmapsto[\widetilde C],
\]
where \(\widetilde C\) is the smooth normalization of \(C\), and \(M_g\) parametrizes smooth connected projective curves of genus \(g\). Let \(Z\) be the reduced closure of its image. Generically finite onto the image means that the induced function-field extension \(k(Z)\hookrightarrow k(V)\) is finite; birational onto the image means that this inclusion is an isomorphism.

\textbf{Question as supplied.} For a general polarized \(K3\) surface \((X,L)\) and an irreducible component \(V\subset V_{n,\delta}(X)\), is \(\psi_{X,n,\delta}|_V\) generically finite, or birational when \(g>1\), onto its image in \(M_g\)? The accompanying assessment describes the intended scope as all components and all admissible \(p,n,\delta\). The source wording does not clarify whether ``or'' distinguishes two proposed assertions or requests generic finiteness together with the stronger birationality assertion for \(g>1\). No unspecified low-genus or numerical hypotheses are added here.
\subsection{Short English statement}
For curves with a fixed number of nodes on a general polarized \(K3\) surface, ask whether the isomorphism class of the smooth normalization determines only finitely many curves in each Severi component, and whether it determines a unique general curve birationally when \(g>1\). The precise parameter ranges, base field, polarization conventions, and relationship between these two assertions require clarification.
\subsection{Sources}
\begin{itemize}
\item Local UnsolvedMath registry, permanent ID \(30002977\), title \emph{Generic Finiteness of K3 Severi Moduli Maps}; exact source problem code \texttt{catalogue entry}. The supplied collection is \emph{Oberwolfach Reports Open Problems}. No individual proposer is supplied; the recorded proposed year is 2015.
\item Supplied report attribution: \emph{Mini-Workshop: Singular Curves on \(K3\) Surfaces and Hyperkähler Manifolds} (2016), Oberwolfach Reports, pp.~2939--2967, DOI \texttt{10.4171/owr/2015/51}. \url{https://doi.org/10.4171/owr/2015/51}. The citation's 2016 date is retained alongside the supplied proposed year 2015 and the DOI containing 2015; their relationship has not been independently checked.
\item Ciliberto--Flamini--Galati--Knutsen, \emph{Moduli of Nodal Curves}. \url{https://doi.org/10.1016/j.aim.2017.01.021}.
\item \emph{The Moduli of Singular Curves on K3 Surfaces}. \url{https://arxiv.org/abs/1401.1047}.
\item \emph{Curves of Maximal Moduli on K3 Surfaces}. \url{https://doi.org/10.1016/j.jms.2021.104832}.
\item Supplied literature assessment, checked \texttt{2026-08-21}: labeled partially solved, with generic finiteness and birationality reported in many ranges, while the question for all components and all \(p,n,\delta\) is reported to remain open. This is supplied dated curation derived from the OWR record, not an independently verified literature assessment; time-sensitive status and cited claims require independent verification.
\item The registry labels the statement \texttt{corrected\_verified} and reports comparison with the official Oberwolfach report. This verification claim is supplied metadata and has not been independently confirmed for this draft.
\end{itemize}
\end{document}
